\section{Scope and conclusion}
\label{sec:conclusion}

The affine coordinate $(p-1)k+1$ links three levels of structure.  Its
power divisibility isolates one hole class modulo $p^N$; its behavior on a
finite interval separates prime from composite bases under the constructive
period convention; and its translates around $r_n$ freeze every bounded
off-center window.  The last fact converts the general finite-local-rule
description into a unique letter quotient.  Once that collapse is known,
the fixed-source factor problem becomes the refinement poset of
independent-set partitions of a finite graph.

All conclusions retain the hypotheses declared in \cref{sec:family}.  The
base is one fixed integer $p\geq3$, the directives are periodic with exact
support and unequal cyclic neighbors, maps are onto and pointed, and both
source and target lie in the affine family.  The target classification is
modulo pointed conjugacy.  We have not classified maps between different
bases, maps over a nonzero odometer element, nonpointed maps, arbitrary
simple Toeplitz systems, or factors outside the family.  The admissible
partition order is a poset; the results do not promote it to a lattice in
general.

Two directions remain natural within these boundaries.  One is to replace
pointedness by a prescribed nonzero odometer phase and determine whether a
translated version of the high-center normal form survives.  The other is
to combine the prime-base constructive obstruction with additional
arithmetic data to study cross-base sufficiency.  Neither question is
settled by the same-base classification proved here.
